\chapter{Methodology}
\label{chap:methodology}

This chapter describes the two experiments conducted in this thesis. It first
summarizes their design (\cref{sec:method-overview}), then presents the
datasets and ground truths used in Experiment~1 (\cref{sec:method-dataset}),
the consolidation pipeline used in Experiment~1 (\cref{sec:method-pipeline}), and the
traceability evaluation used for RQ1 and RQ2
(\cref{sec:method-metrics}). Finally, the full design of Experiment~2, which
addresses semantic conflict handling in RQ3, is presented separately in
\cref{chap:experiment-conflict}.

\section{Overview of the Experiments}
\label{sec:method-overview}

This thesis contains two experiments:

\begin{itemize}[leftmargin=*]
  \item \textbf{Experiment~1: Traceability Evaluation (RQ1 and RQ2).}
  This experiment uses two project histories, three AI coding agents, and four
  consolidation strategies. RQ1 compares the three single-pass strategies:
  \emph{basic}, \emph{loose}, and \emph{strict}. RQ2 evaluates the proposed
  \emph{spec-guided} strategy against these baselines. Together, the four
  strategies produce 24 master-prompt candidates, which are evaluated using
  mean AlignScore and untraceable rate.
  \item \textbf{Experiment~2: Semantic Conflict Handling (RQ3).}
  This experiment begins with the hand-built ground-truth master prompt of each
  project as a settled specification. The simple project receives one
  conflicting instruction and one compatible control, while the hard project
  receives four conflicting instructions and one compatible control. Claude,
  Codex, and Gemini are tested under three conditions:
  \emph{consolidate-then-ask}, \emph{do-not-resolve}, and
  \emph{spec-guided}. The experiment examines how each condition represents
  semantic conflicts without automatically selecting between incompatible
  requirements.
\end{itemize}

The following sections describe the dataset, consolidation pipeline, and
traceability evaluation used in Experiment~1. The complete design and evaluation criteria
of Experiment~2 are presented in
\cref{chap:experiment-conflict}.

\section{Dataset and Ground Truth}
\label{sec:method-dataset}

Two synthetic project histories were created for controlled evaluation,
summarized in \cref{tab:dataset-overview}. They are different in length, complexity,
and writing style.

\begin{table}[htbp]
  \centering
  \caption{Overview of the two evaluation projects.}
  \label{tab:dataset-overview}
  \begin{tabular}{@{}llp{9.7cm}@{}}
    \toprule
    Project & Prompts & Description \\
    \midrule
    simple project & 12 & A small portal project with navigation
      requirements, a conditional color rule, a contact form, and several
      later navigation changes. \\
    hard project & 45 & A larger ERP-style application with customers,
      vendors, products, orders, invoices, reports, roles, and permissions.
      Its history contains more revisions, removals, and conversational
      wording. \\
    \bottomrule
  \end{tabular}
\end{table}

A ground-truth master prompt was constructed manually for each project. The
prompts were read in their original order, and only the requirements that
remained valid at the end of the history were retained. A later modification
replaced the earlier value, while a removal or undo instruction deleted the
affected requirement when no replacement was given. Details that did not
originate from the user's prompts were excluded.

The ground truth was also split manually into atomic requirements. Then, when the
LLM split the candidate master prompts, these ground-truth requirements
were used to check whether a candidate unit combined too many requirements or
split one requirement into unnecessarily small parts.

Every Experiment~1 candidate is evaluated against the ground truth for the
same project. Both ground-truth master prompts are also used as settled
starting artifacts in Experiment~2. The simple-project case introduces one
conflict and one compatible control, while the hard-project case introduces
four conflicts and one compatible control.

\section{Consolidation Pipeline}
\label{sec:method-pipeline}

Each project history was replayed with three AI coding agents:

\begin{itemize}[leftmargin=*]
  \item Claude, using Claude Sonnet 4.6;
  \item Codex, using GPT-5.5-Codex;
  \item Gemini, using Gemini Flash 3.5.
\end{itemize}

Separate sessions were used for each agent, project, and strategy. The project
prompts were delivered in their original order.

Before Prompt~1, every session received the same fixed setup message
(\cref{sec:pre-experiment}). It instructed the agent to treat all next
requests as part of one project and to remain in the same workspace. It also
stated that the setup message was not a software requirement and must not
appear in the master prompt.

\subsection{Consolidation Strategies}

Experiment~1 compares four strategies, summarized in
\cref{tab:consolidation-variants}.

\begin{table}[htbp]
  \centering
  \caption{The four consolidation strategies used in Experiment~1.}
  \label{tab:consolidation-variants}
  \begin{tabular}{@{}lp{9.5cm}@{}}
    \toprule
    Strategy & Description \\
    \midrule
    \texttt{basic\_prompt} & The agent is asked to consolidate the ordered
      prompt history into one master prompt. No explicit traceability rule is
      provided. \\
    \texttt{loose\_consolidate} & The basic instruction is extended by asking
      the agent to reflect later modifications and removals and to avoid
      additional or missing features. \\
    \texttt{strict\_consolidate} & The loose instruction is extended with an
      explicit traceability rule: every requirement must originate from the
      numbered prompts, and unrequested or non-traceable details must be excluded. \\
    \texttt{spec\_guided} & The agent records the raw prompts in
      \texttt{prompts.md} and maintains the current master prompt in
      \texttt{spec.md} after every new instruction. \\
    \bottomrule
  \end{tabular}
\end{table}

The first three strategies are single-pass approaches. The agent receives the
complete history and creates the master prompt once, at the end. Their
instructions are cumulative: \texttt{loose\_consolidate} adds a faithfulness
clause to \texttt{basic\_prompt}, and \texttt{strict\_consolidate} adds an
explicit traceability rule. The exact templates are reproduced in
\cref{app:consolidation-prompts}.

The fourth strategy, \texttt{spec\_guided}, maintains the master prompt
throughout the session under a standing governance instruction. It uses two
files:

\begin{itemize}[leftmargin=*]
  \item \texttt{prompts.md}, an append-only, verbatim log of the numbered
  project prompts;
  \item \texttt{spec.md}, the current master prompt containing only the
  requirements that remain valid.
\end{itemize}

After each new prompt, the agent records it in \texttt{prompts.md} and updates
\texttt{spec.md}. Compatible requirements are added, while clear
modifications, replacements, removals, and undo instructions update the
affected topic immediately.

When a prompt changes an existing topic, the agent reconstructs that topic
from its relevant prompt history instead of editing the current sentence
locally. This helps remove outdated values and conversation-dependent wording.
The development of this rule is described in
\cref{sec:rq2-governance}.

If a new instruction cannot hold together with a settled requirement and does
not clearly indicate a deliberate revision, the agent leaves the current
specification unchanged and records the conflict for the user to decide. The
same conflict-handling principle is applied in Experiment~2 through a
separate governance instruction written for the project-handoff setting.

At the end of the session, \texttt{spec.md} itself is used as the candidate
master prompt. The full governance instruction is reproduced in
\cref{lst:governance}.

This method borrows the general spec-driven development principle that the
specification can act as the main statement of what should be built
\citep{piskala2026specdriven}. It does not implement a complete spec-driven
development toolchain; \texttt{spec.md} is used here only to maintain a
traceable master prompt.

\section{Traceability Metrics}
\label{sec:method-metrics}

Experiment~1 evaluates whether the requirements stated in each candidate are
supported by the ground truth. The evaluation focuses on content that appears
in the candidate but is not part of the final user intent, including
agent-added details and outdated requirements.

\subsection{Atomic Requirement Decomposition}

A whole-document score can hide a small number of unsupported requirements
among many correct ones. Each candidate is therefore evaluated one atomic
requirement at a time.

The decomposition follows the general idea of Claimify
\citep{metropolitansky2025claimify}, but this thesis does not run the Claimify
system. Instead, a fixed prompt designed for software requirements is applied
to each candidate sentence using an LLM
(\cref{sec:decomp-prompt}).

The prompt uses a keep-all approach. It retains unsupported details,
outdated requirements, conversational statements, and any other content in
the candidate so that these units remain available for scoring.

The ground truth and candidates are decomposed differently:

\begin{itemize}[leftmargin=*]
  \item the ground-truth atomic requirements are created manually during
  ground-truth construction;
  \item candidate requirements are first generated by the fixed LLM
  decomposition prompt.
\end{itemize}

The candidate requirements were first split by an LLM using the fixed
decomposition prompt. The LLM did not receive the complete ground-truth master
prompt or its list of atomic requirements for the candidate.
However, the few-shot examples in the prompt followed the same splitting
decisions used when creating the ground truth.

After the LLM produced the candidate requirements, the author compared them
with the manually split ground truth to check that both were divided at a
similar level of detail. If one candidate unit combined two or more requirements 
that were separate in the ground truth, it was divided into smaller units. 
If the LLM split one requirement into unnecessarily small parts, those parts were kept
together. This review changed only how the candidate text was divided; it did
not remove, rewrite, or correct its content.

The ground truth itself was split manually because the decomposition prompt
sometimes divided one enumeration into several weaker requirements. For
example, a requirement defining four navigation buttons could become four
separate requirements, losing the meaning that the four buttons together form
one complete set.

Both the ground-truth construction and the candidate-decomposition review were
performed by the same author. The possible influence of this process on the
evaluation is discussed in \cref{sec:discussion-limitations}.


\subsection{AlignScore}

AlignScore measures whether a candidate requirement is supported by the
ground truth \citep{zha2023alignscore}:

\[
\mathrm{ALIGN}(\text{context}, \text{claim}) \rightarrow [0,1].
\]

For each atomic candidate requirement $r$, the complete ground-truth master
prompt is used as the context:

\[
\mathrm{AlignScore}(r) = \mathrm{ALIGN}(\text{context}=\text{ground truth},\ \text{claim}=r).
\]

The scores are computed in \texttt{nli\_sp} mode using the AlignScore-large
checkpoint with a RoBERTa-large backbone. The comparison in
\cref{sec:demo-experiment} explains why entailment is used instead of word or
embedding similarity.

Two values are reported for each candidate:

\begin{itemize}[leftmargin=*]
  \item \textbf{Mean AlignScore}: the average score across its atomic
  requirements;
  \item \textbf{Untraceable rate}: the percentage of requirements with an
  AlignScore below $0.5$.
\end{itemize}

The untraceable rate indicates how many candidate requirements are not strongly
supported by the ground truth and therefore require further inspection. A
requirement below the threshold is treated as potentially untraceable, not
automatically as invented. As shown in \cref{sec:rq2-results}, some faithful
short requirements also receive unexpectedly low scores.

This is a precision-style evaluation because it checks only the content
included in the candidate. It does not measure whether the candidate omitted
a valid ground-truth requirement. Ground-truth coverage is therefore left as
a limitation (\cref{sec:discussion-limitations}).


